% Lösungen Einheit 1 von 5 – Scharbegriff und Parameterwert (zusammenbau v0.1)
\einheitenkopf{Einheit 1 von 5 · Scharbegriff und Parameterwert}

\erg{10}{a) nach $a$ \quad b) $f_3(x) = 3x^2 - 2x$; $f_3(2) = 8$ \quad c) $f_a(1) = a + a = 10$, also $a = 5$. \quad d) $18e^{-3a} = 18e^{-6}$, also $-3a = -6$ und $a = 2$. \quad e) $p_a(2) = 2$, also $4a = 2$ und $a = 0{,}5$. \quad f) I: $c = 0$, II: $c = 2$, III: $c = -2$ (der y-Achsenabschnitt ist $c$). \quad g) $a = 3$: $f_3(x) = 2x + 3$, Steigung $2$, y-Achsenabschnitt $3$. \quad h) Stützpunkte $(0 | -1)$, $(\pm 1 | 1)$, $(\pm 2 | 7)$; enger als $p_{0{,}5}$. \quad i) Anfangshöhe $H(0) = 10$ cm (1); die Höhe nähert sich langfristig 200 cm, ein Grenzwert, kein Maximum (2); nach 30 Tagen ist die Pflanze 100 cm hoch (3).}

10h)

\begin{ksys}[xmin=-3,xmax=3,ymin=-2,ymax=8,klein] \funktionab{0.5*\x^2-1}{p_{0{,}5}}{-3}{3} \funktionab{2*\x^2-1}{p_2}{-2.1}{2.1} \end{ksys}

\erg{11}{a) Tim hat die Stelle $2$ für $k$ eingesetzt und nach $x$ aufgelöst. Richtig: $g_k(2) = 4k - 3 = 9$, also $k = 3$. \quad b) Mit $a = 2$ entsteht die feste Funktion $f_2(x) = 2x^2 + 1$; jede Stelle hat genau einen Funktionswert, also gibt es genau einen Graphen. \quad c) $b_a(6) = 0$: $36a = 9$, also $a = 0{,}25$; $b(3) = 6{,}75$ m.}
